\normaltrue \difficilefalse \tdifficilefalse
\correctiontrue

\exer{Pompe à palettes$\star$ \label{GEO:03:C2:06:10}}
\setcounter{question}{0}
\marginnote{\xpComp{GEO}{03}}%\UPSTIcompetence[2]{C2-06}
\index{Compétence C2-06}\index{Compétence GEO-03}
\index{Pompe à palettes}

\ifprof
\else
Soit le mécanisme suivant. On a
$\vect{AO}=e\vect{i_0}$ et $\vect{AB}=\lambda(t)\vect{i_1}$.
De plus, $e=\SI{10}{mm}$ et $R=\SI{20}{mm}$.
Le contact entre \textbf{0} et \textbf{2} en $B$ est maintenu en permanence,
notamment par effet centrifuge lors de la rotation de la pompe.
\begin{marginfigure}
\includegraphics[width=\linewidth]{10_01}
\end{marginfigure}
\fi

\question{Tracer le graphe des liaisons.}
\ifprof
\begin{marginfigure}
\includegraphics[width=\linewidth]{10_01_c}
\end{marginfigure}
\fi

\question{Exprimer $\lambda(t)$ en fonction de $\theta(t)$.}
\ifprof
La fermeture géométrique s'écrit
\[
\vect{OA}+\vect{AB}+\vect{BO}=\vect{0}.
\]
En projetant sur $\vect{i_0}$ et $\vect{j_0}$, puis en éliminant
$\varphi(t)$, on obtient
\[
R^2=\left(-e+\lambda(t)\cos\theta(t)\right)^2
+\lambda(t)^2\sin^2\theta(t).
\]
Ainsi,
\[
\lambda(t)^2-2e\lambda(t)\cos\theta(t)+e^2-R^2=0,
\]
d'où
\[
\lambda(t)=e\cos\theta(t)
\pm\sqrt{R^2-e^2\sin^2\theta(t)}.
\]
La solution physique correspond à la racine positive.
\fi

\question{En utilisant Python, tracer $\lambda(t)$ en fonction de $\theta(t)$.}
\ifprof
On conserve la solution positive et on obtient la courbe suivante.
\begin{marginfigure}
\includegraphics[width=\linewidth]{10_02_c}
\end{marginfigure}
\fi

\question{Exprimer $\dot{\lambda}(t)$ en fonction de $\dot{\theta}(t)$.}
\ifprof
En dérivant l'expression précédente, on obtient
\[
\dot{\lambda}(t)
=-e\dot{\theta}(t)\sin\theta(t)
-\frac{e^2\dot{\theta}(t)\sin\theta(t)\cos\theta(t)}
{\sqrt{R^2-e^2\sin^2\theta(t)}}.
\]
\begin{marginfigure}
\includegraphics[width=\linewidth]{10_03_c}
\end{marginfigure}
\fi

\ifprof
\else
On prendra une section de piston \textbf{2} de $\SI{1}{cm^2}$ et une vitesse
de rotation $\dot{\theta}=2\pi\,\SI{}{rad.s^{-1}}$.
\fi

\question{Exprimer le débit instantané de la pompe.}
\ifprof
Le débit instantané associé à un piston est
\[
q(t)=S\dot{\lambda}(t).
\]
\fi

\question{En utilisant Python, tracer le débit instantané de la pompe pour un tour, pour $e=\SI{10}{mm}$ puis $e=\SI{15}{mm}$.}
\ifprof
\begin{marginfigure}
\includegraphics[width=\linewidth]{10_04_c}
\end{marginfigure}
\fi

\question{En utilisant Python, tracer le débit instantané de la pompe pour un tour avec $e=\SI{10}{mm}$ et cinq pistons.}
\ifprof
Les cinq contributions sont décalées de $2\pi/5$.
On annule les valeurs négatives, qui correspondent à la phase d'aspiration,
puis on somme les cinq débits positifs :
\[
q_{\mathrm{tot}}(\theta)
=\sum_{k=0}^{4}
\max\!\left(0,\,q\!\left(\theta-\frac{2k\pi}{5}\right)\right).
\]
La courbe résultante est donnée ci-dessous.
\begin{marginfigure}
\includegraphics[width=\linewidth]{10_05_c}
\end{marginfigure}
\fi

\ifprof
\else
\ifcolle
\else
\marginnote{
\begin{solution}
Indications :
\begin{enumerate}
\item Écrire la fermeture géométrique.
\item Éliminer l'angle $\varphi(t)$ par élévation au carré et sommation.
\item Conserver la racine positive pour $\lambda(t)$.
\item Utiliser $q(t)=S\dot{\lambda}(t)$.
\item Pour cinq pistons, décaler les contributions de $2\pi/5$.
\end{enumerate}
\end{solution}}
\fi
\marginnote{Corrigé voir \ref{GEO:03:C2:06:10}.}
\fi

% BEGIN AUTO CORRIGE INCLUDE
\InclureCorrige{corrige.tex}
% END AUTO CORRIGE INCLUDE
